\section{有限位数境界上の相対 dg smooth proper 族}
\label{sec:dg-smooth-torsion-families}

本節では、構成済みの十七次元 Jacobian atlas の中で、
十六次元の局所閉部分族の非空開上に相対 dg smooth proper 性を証明する。
このような部分族の合併は atlas に稠密である。
十七次元 atlas 自体の非空開は、この結果には含めない。
五線 compact generator の係数上 perfect 性を既存の入力とし、
有限中心の射影模型に対する大域対角の議論を加える。

\subsection{係数上 perfect な dg 代数の一般点からの広がり}
\label{sec:dg-generic-spreading}

以下では dg 代数は結合的かつ単位的であり、テンソル積と enveloping algebra は
すべて導来の意味で用いる。\(R\) 上の dg 代数 \(E\) に対し
\(E^e=E^{\mathrm{op}}\otimes_R^{\mathbf L}E\) と書く。
\(E\) が \(R\)-proper であるとは \(E\in\operatorname{Perf}(R)\)、
\(R\)-smooth であるとは対角加群 \(E\in\operatorname{Perf}(E^e)\)
であることをいう。

\begin{lemma}[一般点からの smoothness の広がり]
\label{lem:dg-generic-spreading}
\(R\) を整域である Noether 環、\(K=\operatorname{Frac}(R)\) とし、
\(E\) を \(E\in\operatorname{Perf}(R)\) を満たす \(R\)-dg 代数とする。
\(E_K=E\otimes_R^{\mathbf L}K\) が \(K\)-smooth ならば、
ある \(0\ne f\in R\) に対し \(E_f\) は \(R_f\)-smooth proper である。
したがって任意の \(R_f\)-代数 \(S\) に対し \(E_f\otimes_{R_f}^{\mathbf L}S\)
は \(S\)-smooth proper である。
\end{lemma}

\begin{proof}
\(D=E^e\) とおく。\(E\) が係数上 perfect なので \(D\) も
\(R\)-perfect である。\(R\) の非零元による局所化の filtered system に対し、
perfect dg 加群は有限段階で定義できる：
\[
 \mathop{\mathrm{colim}}_{0\ne g\in R}\operatorname{Perf}(D_g)
 \simeq \operatorname{Perf}(D_K).
\]
これは To\"en--Vaqui\'e\cite[Lemma~2.10]{ToenVaquieObjects}、または
To\"en\cite[事実 F4]{ToenDefinitionRings}による。この連続性に \(D\) の homotopical finite presentation
は必要ない。有限 cell 加群とその retract の有限個のデータを降下させる性質である。

仮定より \(E_K\in\operatorname{Perf}(D_K)\) であるから、
ある \(g\ne0\) と \(P\in\operatorname{Perf}(D_g)\) を選んで
\(P_K\simeq E_K\) とできる。この同型を \(\alpha_K\) とする。
source \(P\) の perfectness により、任意の \(D_g\)-加群 \(M\) について
\[
 \mathbf R\operatorname{Hom}_{D_g}(P,M)\otimes_{R_g}^{\mathbf L}K
 \simeq
 \mathbf R\operatorname{Hom}_{D_K}(P_K,M_K)
\]
が成り立つ。自由加群 \(D_g[n]\) に対して明らかであり、
有限直和・有限 cone・retract に関して保たれるためである。
ここでは target \(M\) の \(D_g\)-perfectness を仮定していない。
したがって \(M=E_g\) として \(\alpha_K\) の有限個の分母を払い、
さらに局所化し、記号を変えずに書けば射
\(\alpha:P\longrightarrow E_g\) を得る。

その cone \(C\) は \(R_g\)-perfect である。
実際、\(D_g\) は \(R_g\)-perfect なので制限関手は
\(\operatorname{Perf}(D_g)\) を \(\operatorname{Perf}(R_g)\) に送る。
よって \(P\) と \(E_g\) はともに係数上 perfect である。
また \(C_K=0\) である。
係数上 perfect な複体のコホモロジーは有限個の有限生成加群であり、
各 \(H^i(C)\) は \(K\) への局所化後に消える。
有限個の生成元を殺す非零元を掛け合わせると、さらにある局所化上で
\(C=0\) となる。したがってその上で \(E\simeq P\) は
\(E^e\)-perfect、すなわち smooth である。
properness および smoothness の任意の導来基底変換による保存は、
perfect 複体および perfect 対角加群の基底変換から従う。
\end{proof}

\begin{lemma}[体拡大からの降下]
\label{lem:dg-smooth-field-descent}
\(k\subset L\) を体拡大、\(E\) を \(k\)-proper な dg 代数とする。
\(E_L\) が \(L\)-smooth ならば \(E\) は \(k\)-smooth である。
\end{lemma}

\begin{proof}
\(D=E^{\mathrm{op}}\otimes_k E\) とする。
\(L\) をその有限生成 \(k\)-部分代数の filtered colimit とみなし、
上記と同じ perfect 加群の連続性を適用する。
ある有限生成整域 \(A\subset L\) と \(P\in\operatorname{Perf}(D_A)\)
および同型 \(P_L\simeq E_L\) が存在する。
source \(P\) の compactness によって、\(A\) を必要なら大きくして、
この同型を射 \(\alpha:P\to E_A\) として定義できる。
cone \(C\) は \(A\)-perfect であり \(C\otimes_A L=0\) を満たす。
\(\operatorname{Frac}(A)\subset L\) の忠実平坦性から
\(C_{\operatorname{Frac}(A)}=0\)、したがって非零元を一つ逆にすると
\(C=0\) とできる。こうして \(E_A\in\operatorname{Perf}(D_A)\)
となる有限生成 \(k\)-整域 \(A\subset L\) が得られた。

\(\operatorname{Spec}(A)\) の閉点を取り、その剰余体を \(k'\) とする。
Zariski の補題により \([k':k]<\infty\) である。
\(E_{k'}\) は \(D_{k'}\)-perfect であり、
\(D_{k'}\) は \(D\)-加群として有限自由なので、
制限した \(E_{k'}\) は \(D\)-perfect である。
\(\lambda(1)=1\) を満たす \(k\)-線型写像 \(\lambda:k'\to k\) を選ぶと
\[
 E\longrightarrow E\otimes_k k'
   \xrightarrow{\,1\otimes\lambda\,}E
\]
の合成は恒等写像である。いずれも \(D\)-線型なので、
\(E\) は \(D\)-perfect 加群の retract、したがって \(D\)-perfect である。
\end{proof}

\subsection{大域対角を扱うための補題}

\begin{proposition}[左右の正則性から大域 dg smooth 性へ]
\label{prop:seed-finite-central-global-smooth}
$k$ を完全体、$Y$ を有限型分離 $k$-スキーム、$\mathscr A$ を
coherent な $\mathcal O_Y$-代数とする。
左右それぞれについて
\[
 \mathrm D^b(\operatorname{coh}\mathscr A)
   =\operatorname{Perf}(Y,\mathscr A),\qquad
 \mathrm D^b(\operatorname{coh}\mathscr A^{\mathrm{op}})
   =\operatorname{Perf}(Y,\mathscr A^{\mathrm{op}})
\]
と仮定する。すると $\operatorname{Perf}_{\mathrm{dg}}(Y,\mathscr A)$ は
$k$ 上 dg smooth である。さらに $Y$ が $k$ 上 proper なら、この dg 圏は proper である。
\end{proposition}

\begin{proof}
$Y$ は noetherian J-2 スキームである。
Elagin--Lunts--Schn\"urer\cite[Theorem~4.15]{ELSSmoothAlgebras}により
$\mathrm D^b(\operatorname{coh}\mathscr A)$ は classical generator $G$ を持つ。
仮定により $G$ は $\mathscr A$-perfect である。
以下、$\mathscr A$-加群は右加群とする。
双対
\[
 G^\vee=\mathbf R\mathcal Hom_{\mathscr A}(G,\mathscr A)
\]
は右 $\mathscr A^{\mathrm{op}}$-perfect であり、perfect 加群の双対は反同値なので、
第二の仮定により $G^\vee$ は
$\mathrm D^b(\operatorname{coh}\mathscr A^{\mathrm{op}})$ の classical generator となる。

$Y\times_kY$ 上で
\[
 \mathscr B=\mathscr A^{\mathrm{op}}\boxtimes_k\mathscr A,
 \qquad K=G^\vee\boxtimes_kG
\]
と置く。Elagin--Lunts--Schn\"urer\cite[Theorem~4.18]{ELSSmoothAlgebras}は、完全体上の
noetherian J-2 スキームと coherent 代数に対して、
\emph{任意に選んだ}二つの classical generator の外積が
外積代数の有界 coherent 導来圏を厚く生成することを述べている。
従って
\[
 \mathrm D^b(\operatorname{coh}\mathscr B)=\operatorname{thick}(K).
\]
分離性により対角 $\delta:Y\hookrightarrow Y\times_kY$ は閉埋込みであり、
$\delta_*\mathscr A$ は coherent な右 $\mathscr B$-加群である。
右作用は、局所的に
$m\cdot(a^{\mathrm{op}}\otimes b)=amb$ と与える。
ゆえに $\delta_*\mathscr A\in\operatorname{thick}(K)$ である。

$E=\mathbf R\operatorname{Hom}_{\mathscr A}(G,G)$ と置く。
有限 affine 被覆の \v Cech 模型と perfect 加群の tensor--Hom 恒等式から
\[
 \mathbf R\operatorname{Hom}_{\mathscr B}(K,K)
       \simeq E^{\mathrm{op}}\otimes_k E,
 \qquad
 \mathbf R\operatorname{Hom}_{\mathscr B}(K,\delta_*\mathscr A)
       \simeq E
\]
を得る。後者は $E^{\mathrm{op}}\otimes_kE$-加群の同型であり、右辺は対角作用を持つ。
次数零では $f\cdot(e_1^{\mathrm{op}}\otimes e_2)=e_1fe_2$ であり、
一般の次数では通常の Koszul 符号を入れる。
この比較は局所的に $G$ を有限生成射影 $\mathscr A$-加群の有界複体で表して確かめ、
\v Cech 複体で大域化できる。
$\mathscr A$ が $\mathcal O_Y$ 上平坦であるとは仮定していない。
外積の基底が体 $k$ であることと、$\mathscr A$-perfect 性を用いる。

第一の同型により、$\mathbf R\operatorname{Hom}_{\mathscr B}(K,-)$ は
$\operatorname{thick}(K)$ を
$\operatorname{Perf}(E^{\mathrm{op}}\otimes_kE)$ に移す。
第二の同型と $\delta_*\mathscr A\in\operatorname{thick}(K)$ から、
対角 $E$ は $E^{\mathrm{op}}\otimes_kE$ 上 perfect である。
Morita 不変性により、これが所望の dg smooth 性である。

$Y$ が proper なら、二つの perfect $\mathscr A$-複体の導来内部 Hom は
有界 coherent $\mathcal O_Y$-複体である。
その大域切断は有界かつ各次数有限次元なので、proper 性も従う。
\end{proof}

\begin{remark}[局所対角と大域対角の役割]
有限中心の affine chart の対角が perfect であることだけを、
大域 dg smooth 性の証明と同一視していない。
上の証明で大域化を担う入力は、外積 classical generator の生成定理と
$\delta_*\mathscr A$ に対する dg Morita 比較である。
この命題の一般形の初出は未確認であり、
主張の新規性は既知の Elagin--Lunts--Schn\"urer の生成定理と区別する。
\end{remark}

\subsection{有限位数境界の点は二係数開集合に稠密である}

$T_0=\operatorname{Spec}\mathbb C[a,b,\Delta^{-1}]$ とし、
現行原稿の規約
\[
 p=2-b,\qquad q=-(2+b),\qquad r=-6a
\]
を用いる。$p-q=4$ なので
\[
 \mathbb A^2_{a,b}\xrightarrow{\sim}
 \{[p:q:r]\in\mathbb P^2\mid p-q\ne0\},
 \qquad
 a=-\frac{2r}{3(p-q)},\quad
 b=-\frac{2(p+q)}{p-q}.
\]
従って $T_0$ はこの射影平面の非空開集合と同一視できる。

\begin{lemma}[torsion 境界パラメータの稠密性]
\label{lem:seed-torsion-boundary-dense}
$T_0(\mathbb C)$ において、楕円境界の自己同型 $\sigma$ が
有限位数を持つ点の集合は Zariski 稠密である。
\end{lemma}

\begin{proof}
$pqr$ が可逆な開集合上で
\[
 \mu=\frac{p^3+q^3+r^3}{pqr}
\]
を取る。$\Delta\ne0$ の条件は $\mu^3\ne27$ を含む。
$\mu$ のファイバーは、滑らかな Hesse 三次曲線
\[
 E_\mu:\quad X^3+Y^3+Z^3-\mu XYZ=0
\]
の開部分である。
Sklyanin の同定により、$O=[1:-1:0]$ を原点とすると、
そのパラメータにおける $\sigma$ は
$P=[p:q:r]\in E_\mu$ による平行移動である。
ここで関係の規約は
$pyz+qzy+rx^2=0$ とその巡回式であり、
Walton--Wang--Yakimov\cite[Definition--Lemma~2.13]{WWYSklyaninPoisson}と一致する。

任意の非空開集合 $W\subset T_0$ を取る。
ある閉点 $w\in W$ と $\mu_0=\mu(w)$ を選ぶと、
$W\cap E_{\mu_0}$ は滑らかな既約楕円曲線 $E_{\mu_0}$ の非空開集合である。
複素楕円曲線の torsion 点は無限集合であり、
既約曲線の真の閉集合は有限なので Zariski 稠密である。
従って $W\cap E_{\mu_0}$ に torsion 点 $P$ が存在する。
その平行移動 $\sigma$ は有限位数である。
これで任意の非空 $W$ との交わりが得られた。
\end{proof}

\begin{remark}
torsion の位数を一つに固定していない。
任意の指定された有限個の位数を除いても、
各楕円曲線上には無限個の torsion 点が残るので同じ証明が成立する。
特に $\Delta$ の小位数除外条件は障害にならない。
一方、位数6の既存例 $B_*$ 自身を使うことはできない。
その最高次パラメータは $a=1/3,b=0$、従って
$(p,q,r)=(2,-2,-2)$ であり、$q^3-r^3=0$ のため $T_0$ の外にある。
\end{remark}

\subsection{固定した有限位数 ambient の相対開}

\begin{proposition}[有限位数固定 ambient 上の相対 smooth proper 開]
\label{prop:seed-fixed-ambient-relative-smooth}
$t\in T_B(\C)$ を境界 translation が有限位数となる固定 ambient 点とし、
$A_{t,g}=B_t/(\mathcal G_t+\sum_{j=0}^3g_ju^{3-j}v^j)$ と置く。
$P=\mathbb C[g_0,g_1,g_2,g_3]$ とする。
ある非零元 $h\in P$ が存在し、任意の可換 $P_h$-代数 $S$ に対して、
\[
 \mathrm D(\operatorname{QGr}A_{t,S})^c
\]
は $S$ 上 dg smooth proper となる。
この同定と dg Morita 模型は任意の基底変更と両立する。

\end{proposition}

\begin{proof}
系\ref{cor:relative-cubic-compact-proper}と
命題\ref{prop:general17-relative-cy-koszul}により、
\[
 T=\bigoplus_{i=0}^{4}\mathcal O(i),\qquad
 E=\mathbf R\operatorname{Hom}(T,T)
\]
を取ると、$T$ は compact generator、$E$ は $P$-perfect であり、
全ての可換 $P$-代数 $S$ について
\[
 E\otimes_P^{\mathbf L}S
   \simeq\mathbf R\operatorname{Hom}(T_S,T_S)
\]
が成り立つ。この比較は乗法を保つ canonical な dg 比較である。

$K=\operatorname{Frac}P$ とする。
固定した有限位数境界についての全四係数族は、
定理\ref{thm:central-model-finite} により
有限型次数付き可換 $P$-代数 $C$ 上有限である。
従って $Y=\operatorname{Proj}C^{(\ell)}$ と
$\mathscr A=\widetilde{A^{(\ell)}}$ は
$K$ と $\overline K$ への基底変更後にも使える。
$C_S\to\mathscr A_S$ の単射性や全中心の基底変更可換性は不要であり、
中心作用と有限加群性だけを用いる。

定理\ref{CSBGthm:global-open}を左右に適用して得た二つの非零多項式は、
generic point $\operatorname{Spec}\overline K\to\mathbb A^4$ ではともに可逆である。
同定理の幾何学的基底変更版から、
$A_{t,\overline K}$ の左右 tails 大域次元はともに3である。
左右の tails 大域次元が有限であるから、
補題\ref{CSBGlem:global-finite-to-perfect}により有界 coherent 導来圏は
全 twist で厚く生成される。各 twist は局所 $\mathscr A$-perfect なので、
左右とも $\mathrm D^b(\operatorname{coh}\mathscr A)=\operatorname{Perf}(Y,\mathscr A)$ となる。
命題\ref{prop:seed-finite-central-global-smooth}を
$(Y_{\overline K},\mathscr A_{\overline K})$ に適用して
$E_{\overline K}$ の dg smooth 性を得る。
補題\ref{lem:dg-smooth-field-descent}で $E_K$ へ降下する。

補題\ref{lem:dg-generic-spreading}を整域 $P$ と $P$-perfect dg 代数 $E$ に適用すれば、
ある $0\ne h\in P$ について
\[
 E_h\in\operatorname{Perf}
    (E_h^{\mathrm{op}}\otimes_{P_h}^{\mathbf L}E_h)
\]
となる。任意の $P_h\to S$ の導来基底変更はこの perfect 対角を保つ。
$E_h$ 自身の $P_h$-perfect 性も基底変更で保たれるので、
$E\otimes_P^{\mathbf L}S$ は $S$ 上 smooth proper である。
普遍 generator とその endomorphism の基底変更比較により、
これを $\mathrm D(\operatorname{QGr}A_{t,S})^c$ へ移す。
\end{proof}

\begin{remark}
$t\in T_B^e$ の場合、この相対四次元開は新しい十七次元 atlas の中で
ambient $t$ を固定した fiber 内の非空開集合である。
十七次元 atlas 自体の非空開集合とは区別する。
\end{remark}

\subsection{十五次元および十六次元の相対 smooth proper 族}

\begin{proposition}[有限位数境界を固定した全混合変形]
\label{prop:seed-torsion-base-relative-smooth}
$Z\subset T_0$ を integral affine な局所閉部分スキームとし、
その幾何学的一般点における境界 translation が有限位数を持つと仮定する。
十一個の mixed 係数と四個の cubic 係数を全て動かす族
\[
 X_Z=Z\times\mathbb A^{11}_{\alpha,\beta,m_{ij}}
           \times\mathbb A^4_{g_0,g_1,g_2,g_3}
\]
には、非空主開集合 $U_Z$ があり、その上の compact tails は
相対 dg smooth proper である。
この主張は $U_Z$ からの任意の可換係数環への基底変更に両立する。
従って $Z$ が有限位数境界の複素点なら、
十五次元 affine 空間の非空開上の相対 dg smooth proper 族が得られる。
\end{proposition}

\begin{proof}
$R_0=\Gamma(Z,\mathcal O_Z)$、
$R=R_0[\alpha,\beta,m_{ij},g_0,g_1,g_2,g_3]$ と置く。
この $R$ は有限型 $\mathbb C$-整域である。
普遍五線 generator の dg 自己準同型代数 $E_R$ は $R$-perfect であり、
任意の基底変更と両立する。
\[
 F=\operatorname{Frac}(R_0)(\alpha,\beta,m_{ij}),\qquad
 K=F(g_0,g_1,g_2,g_3)=\operatorname{Frac}(R)
\]
と置き、$\overline K$ の中で $F$ の代数閉包 $\overline F$ を選ぶ。
$\overline F$ 上の ambient は固定され、
その境界 translation は仮定により有限位数である。

定理\ref{CSBGthm:global-open}は任意の代数閉標数零体上で成立する。
実際、その有限中心入力は Rains\cite[Theorem~2.13]{Rains}、
開性入力は Rennemo--Segal--Van den Bergh\cite[Proposition~2.1]{RSVdB}であり、
両者ともこの体の仮定で成立する。
二つの中心 chart、行列因子化、境界の twists 生成、Serre 双対性の議論も
\(\mathbb C\) に固有の位相的性質を使わない。
従って $B_{\overline F}$ とその反対環にそれぞれ適用すると、
四係数空間上の二つの非零正則性多項式を得る。

$g_0,g_1,g_2,g_3$ は $\overline F$ 上も代数的独立なので、
この二つの多項式は幾何学的一般点で消えない。
従って $A_{\overline K}$ の左右 tails 大域次元は3である。
被約 Jacobson 係数環の交換子行列を使う
補題\ref{CSBGlem:global-center-lifting}と定理\ref{thm:central-model-finite}の証明も
この体上で成立する。従って有限中心射影模型が存在する。
前命題と同じ左右の twist 生成により、その有界 coherent 導来圏は局所 perfect 圏に一致するので、
命題\ref{prop:seed-finite-central-global-smooth}から
$E_{\overline K}$ は smooth である。
補題\ref{lem:dg-smooth-field-descent}で $E_K$ へ降下し、
補題\ref{lem:dg-generic-spreading}で $R_h$ 上へ延ばす。
$U_Z=D(h)$ とすれば、相対 smooth proper 性と任意の底変更への両立が従う。
全 $X_Z$ 上で中心を一様に選ぶことは、この証明の仮定ではない。
\end{proof}

\begin{theorem}[十六次元 smooth proper 族の稠密な集合]
\label{thm:seed-sixteen-dimensional-dense-families}
十七次元 atlas $T_B^e\times\mathbb A^4$ の任意の非空開集合 $W$ に対し、
smooth integral affine 曲線 $Z\subset T_0$ と、
\[
 V\subset (Z\times\mathbb A^{15})\cap W
\]
なる非空開部分スキームが存在して、
$V$ 上の compact tails は相対 dg smooth proper となる。
$V$ は十六次元であり、任意の可換係数環への基底変更に両立する。
従ってこの atlas では、十六次元の相対 smooth proper 族の像の合併が稠密である。
\end{theorem}

\begin{proof}
Hesse elliptic scheme を
\[
 \mathcal E\longrightarrow M
   =\operatorname{Spec}\mathbb C[\mu,(\mu^3-27)^{-1}],
 \qquad X^3+Y^3+Z^3-\mu XYZ=0,
 \qquad O=[1:-1:0]
\]
とする。
$pqr\ne0$ 上では $\mu=(p^3+q^3+r^3)/(pqr)$ によって
$\mu$ が一意に決まるので、
$T_0$ は $\mathcal E$ の非空開部分スキームと同一視できる。
正整数 $n$ に対し、乗法写像 $[n]:\mathcal E\to\mathcal E$ の
相対微分は $n$ 倍なので可逆であり、滑らかな $M$-曲線間の Jacobian 判定から
この射は \emph{\'etale} である。各幾何学的ファイバーの乗法写像は有限である
\cite[Lemmas~39.9.8,39.9.9]{StacksEllipticTorsion}。
従って $[n]$ は proper かつ quasi-finite、ゆえに有限であり、
その kernel $\mathcal E[n]$ は $M$ 上有限 \emph{\'etale} である。
従って $\mathcal E[n]\cap T_0$ の各非空連結成分は
smooth integral affine 曲線である。
affine 性は、$\mathcal E[n]\cap T_0$ が affine な $T_0$ の閉部分スキームであり、
smooth 曲線の連結成分が開かつ閉であることから従う。

$W$ の $T_0$ への射影像は非空開集合である。
補題\ref{lem:seed-torsion-boundary-dense}により、そこに torsion 点 $P$ を選べる。
$nP=0$ となる $n$ を取り、
$P$ を含む $\mathcal E[n]\cap T_0$ の連結成分を $Z$ とする。
その全ての点で $\sigma^n=1$ なので、
$Z$ は命題\ref{prop:seed-torsion-base-relative-smooth}の仮定を満たす。

$P$ は $W$ の射影像に入るので、
$W\cap(Z\times\mathbb A^{15})$ は空でない。
これは integral な十六次元スキーム $Z\times\mathbb A^{15}$ の開部分である。
命題\ref{prop:seed-torsion-base-relative-smooth}が与える非空主開 $U_Z$ と交わらせ、
\[
 V=U_Z\cap W
\]
とすれば、非空性、次元、相対 smooth proper 性を得る。
\end{proof}

\begin{remark}[十七次元の開性との区別]
前定理では、各 $Z$ 上で有限位数の上界 $n$ を固定している。
得られる十六次元族の合併は稠密であるが、
十七次元 atlas の generic point を含むとは述べていない。
十七次元の相対 dg smooth proper 非空開の構成には、
一般の無限位数境界を扱う別の議論が必要である。
これは第\ref{sec:kuznetsov-moduli-morphism}節で残余圏の幾何性から実現する。
\end{remark}


\begin{corollary}[smooth proper 点の稠密性]
\label{thm:seed-smooth-points-dense}
$T_B^e\times\A^4$ および $\mathfrak C_{\mathrm J}$ の中で、
compact tails が dg smooth proper となる複素点の集合は Zariski 稠密である。
\end{corollary}
\begin{proof}
前定理で得た非空な有限型 $\C$-スキーム $V$ は複素点を持つ。
従って atlas の任意の非空開集合はそのような点を含む。
atlas の全射性から、$\mathfrak C_{\mathrm J}$ の非空開部分の逆像も非空であり、
同じ主張が stack 上で従う。
\end{proof}

\begin{remark}[結論の範囲と既知入力]
外積 generator の定理、perfect dg 加群の filtered colimit、
楕円曲線の torsion と既存の有限中心 Bertini は既知の入力である。
本節の具体的な Jacobian 族への適用と、その相対十六次元族の初出は未確認である。
$V$ は atlas 内の局所閉部分スキームであり、その像が独立に閉部分 stack を
定めることや、異なる torsion 曲線の像が互いに異なることまでは主張しない。
十七次元の一般点の dg smooth 性とその非空開は
定理\ref{thm:kuznetsov-cubic-moduli-morphism}で別に証明する。
全 smooth locus の開性と、全 Abel モジュライ内での成分性は依然として主張しない。
\end{remark}
